% Paste-ready ICCD text for the missing baseline-vs-SPICE stress comparison.
% Requires \usepackage{booktabs}; use table* in two-column IEEE format.

\begin{table*}[t]
\centering
\caption{High-utilization Top-$K$ stress study with all baselines. The setting
uses 768 decoding steps, 64 experts, 8\,MB expert transfers, and the same tight
cache budget for every policy. Lower TPOT and fallback are better.}
\label{tab:supp_topk_baselines}
\scriptsize
\begin{tabular}{rrrrrrr}
\toprule
Top-$K$ & Naive & LRU & MoE-Offload & Pre-gated & SPICE & SPICE fallback (\%) \\
\midrule
2 & 50.42 & 42.92 & 42.92 & 41.66 & \textbf{41.01} & \textbf{0.46} \\
4 & 60.83 & 45.91 & 45.91 & 43.30 & \textbf{42.51} & \textbf{7.43} \\
6 & 71.25 & 49.10 & 49.10 & 45.09 & \textbf{43.14} & \textbf{6.97} \\
8 & 81.67 & 52.65 & 52.65 & 54.37 & \textbf{51.11} & \textbf{24.36} \\
10 & 92.08 & \textbf{56.25} & \textbf{56.25} & 68.90 & 58.36 & 33.41 \\
12 & 102.50 & \textbf{58.39} & \textbf{58.39} & 76.54 & 65.35 & 39.03 \\
\bottomrule
\end{tabular}
\end{table*}

\textbf{High-utilization baseline comparison.}
Table~\ref{tab:supp_topk_baselines} directly addresses the case where baseline
PCIe utilization is already high: the Naive baseline drives modeled PCIe
activity from 20.66\% at $K{=}2$ to 60.98\% at $K{=}12$, while Pre-gated reaches
34.91--100\% activity for $K{\ge}4$. SPICE remains fastest through $K{=}8$ by
reducing synchronous fallback traffic, but its advantage disappears once dense
routing pushes transfer demand beyond the overlap window. At $K{=}10$ and
$K{=}12$, the conservative LRU/MoE-Offloading cache has lower TPOT than SPICE
because additional speculation consumes bandwidth that cannot be hidden. We
therefore state SPICE as a bounded systems optimization: it improves offloaded
MoE serving when useful copy/compute overlap windows exist, but it does not
remove the physical host-to-device bandwidth limit.
